\documentclass{article}

\newcommand{\doctitle}{Assignment 4} % change document title
\usepackage{ellis_header}

\begin{document}

\section{Chapter 2.1}
\setcounter{problem}{3}
\begin{problem} % Problem 4
Consider the possible functions $f: [7] \rightarrow [9]$
\begin{enumerate}[(a)]
    \item How many have $f(3) = 8$? 
    \item How many have $f(3) \neq 8$?
    \item How many have $f(1) \neq 5$ and are one-to-one?
    \item How many have $f(i)$ even, for all $i$?
    \item How many have $\text{rng}(f) = \{5, 6 \}$?
    \item How many in which $f^{-1}$ is not a function?
\end{enumerate}
\end{problem}

\textbf{Answer:}
\begin{enumerate}[(a)]
    \item $f(3)$ has 1 choice. This leaves 6 other values that can go into 9 other bins. I have nine bins, that I can place six values in, with repitition allowed. This is $9^6$.
    \item $f(3)$ has 8 choices, while the six other values have 9 choices. This is $8 \cdot 9^6$.
    \item $f(1)$ has 8 choices, $f(2)$ has eight choices (since one was removed by $f(1)$.) Then $f(3)$ has seven choices, $f(4)$ has six choices... This is $8 \cdot (8)_6 = \dfrac{8 \cdot 8!}{(8-6)!}$.
    \item Instead of 9 choices, we now have 4 choices. So 7 values have 4 choices. This is $4^7$.
    \item Now we have only two choices, $\{5, 6\}$ for 7 values. This is $2^7$. But, we must guarentee at least one values in $\{5\}$ and one value in $\{6\}$. So subtract out when all the values map to 5 or all the values map to 6, which is two cases. Final answer is $2^7 - 2$.
    \item For $f^{-1}$ to be a function, the original function $f$ must be one-to-one. We are looking for where the original function is not one-to-one. Base options is 9 choices for 7 values. So $9^7$. But, we need to subtract out the one-to-one option, which is $(9)_7$. The final answer is $9^7 - (9)_7$.
\end{enumerate}

\setcounter{problem}{6}
\begin{problem} % Problem 7
Find the number of onto functions $[k] \rightarrow [4]$.
\end{problem}

\textbf{Answer:} I am assuming that $k \ge 4$. The base case is $4^k$ options for distribution. Then need to subtract out the following:
\begin{enumerate}[(a)]
    \item Missing 3 of the elements in $[4]$. Answer is $4^1 = 4$.
    \item Missing 2 of the elements in $[4]$. You have $\binom{4}{2} = 6$ elements to miss. You then have $2^k$ ways to distribute the values to the remaining elements. Each distribution has 2 ways that all the values end up on one elements instead of both, so remove those. The answer is $6(2^k - 2)$.
    \item Missing 1 of the elements in $[4]$. You have $\binom{4}{3}= 4$ different elements to miss. Once you miss, you can select $3^k$ places to go. But you need to guarentee that you don't miss more than 1 element. There are 3 ways to miss 2 elements with $2^k$ places to go. There are 3 ways to miss 3 elements with $1^k$ places to go. The answer is $4(3^k - 3 (2^k - 2) - 3 \cdot 1^k)$
\end{enumerate}

Final answer is $4^k - (4 + 6(2^k -2) + 4(3^k -(3 + 3( 2^k - 2))))$

\setcounter{problem}{10}
\begin{problem} % Problem 11
Give a combinatorial proof. For $n \ge 1$ and $k \ge 1$, then $2^{kn} > \max \{n^k, k^n\}$. Hint: Compare relations to functions.
\end{problem}

\textbf{Answer:} Before diving into the details with the general case, I'm going to do a numerical example to get us warmed up. Let us define two sets.
\begin{equation*}
    K = [4] = \{1, 2, 3, 4\} \quad N = [3] = \{1, 2, 3\}
\end{equation*}

All the possible relations from $K$ to $N$ are contained in the Cartesian product of $K \times N$.
\begin{equation}
    K \times N = \{\{1,1\}, \{1, 2\}, \{1, 3\} \dots \{4, 1\}, \{4, 2\}, \{4, 3\}\}
\end{equation}

Every possible relation on $K \rightarrow N$ is a subset of $K \times N$. For instance, for the relation $>$ would have $\{\{2, 1\}, \{3, 1\}, \{3, 2\}, \{4, 1\}, \{4, 2\}, \{4, 3\}\}$. So all the possible relations is $\mathscr{P}(K \times N)$, the power set of the cartesian product of $K \times N$. So, the number of possible relations is $2^{|K \times N|}$. If $k = |K|$ and $n = |N|$ then it is $2^{kn}$.

\medskip
Now, we focus our attention to functions instead of relations. The basic defintion of a function is that every input can only have one output. Thinking about this as distributions, the first number in $K$ has 3 potential choices. The second number in $K$ also has 3 choices, and so one. As a result, when looking at functions, there are $n^k = 3^4 = 81$ possible functions from $K \rightarrow N$. Likewise, if we go the opposite way, from $N \rightarrow K$, there are $k^n = 4^3 = 64$ possible functions.

\medskip
What this is saying, is the number of possible relations between two sets will always be bigger than the possible functions from one set to another and vice a versa.

\begin{proof}
Suppose that $k, n \ge 1$ and $K,N$ are sets of the first postive integers $k,n$ denoted as $K = [k]$ and $N = [n]$. Then the set of all possible relations between $K$ and $N$ is the powerset of $K \times N$, denoted as $\mathscr{P}(K \times N)$. The cardinality of the powerset is $|\mathscr{P}(K \times N)| = 2^{kn}$, or the number of possible relations.

\medskip
For a relation to be considered a function, every value in the domain must map to only one value in the codomain. The function $f: K \rightarrow N$ has $n^k$ possible variations. The set of all function variations is a subset of all the relation variations. As such, $n^k \le 2^{nk}$. 

\medskip
The inequality can be written $n^k < 2^{nk}$ by recognizing the empty set is inside the powerset but will never be considered to meet the definition of a function.

\medskip
Now reversing the roles of $K$ and $N$ with the function $f: N \rightarrow K$ which has $k^n$ possible variations. The same argument from above is used to prove $k^n < 2^{nk}$. This proves that $2^{kn} > \max \{n^k, k^n\}$.
\end{proof}

\section{Chapter 2.2}
\setcounter{problem}{3}
\begin{problem} % Problem 4
Give combinatorial or bijective proofs of the following. Part of your job is to determine all values of $n, k, \text{and / or } m$ for which the identities are valid.
\begin{enumerate}[(a)]
    \item $3^n = \sum_{k=0}^{n} \binom{n}{k} 2^{n-k}$
    \item $\binom{n}{k} \binom{k}{j} = \binom{n}{j} \binom{n-j}{k-j}$
    \item $\multiset{n}{k} = \multiset{k+1}{n-1}$
\end{enumerate}
\end{problem}

\textbf{Answer:}
\begin{enumerate}[(a)]
    \item LHS: How many distributions can be made with n distinct objects (different colored balls) to 3 distinct recipients, $\{\text{Moe, Curly, Larry}\}$.
    
    \medskip
    RHS: Focus attention on the first receipient, Moe. Start with $k=0$, meaning Moe will get all the combinations of zero balls. The remaining balls $n-k$ are distributed in various manners to Curly and Larry. Next, $k=1$, meaning Moe will get all the combinations of one distinct ball. Curly and Larry receive the number of balls that Moe did not get. This pattern is repeated until Moe gets the combinations of all the balls and Curly / Larry receive none. These different cases are added together.

    \medskip
    Example is made with two balls, red and blue $\{R, B\}$. The table below shows $k = 0,1,2$ and the various ball distributions. As known, $3^2 = 9$ which is the number of permutations below.

    \begin{center}
    \begin{tabular}{ |c|c|c|c| } 
    \hline
    k & Moe & Curly & Larry \\ 
    \hline
    0 & $\emptyset$ & R,B&  $\emptyset$ \\
    0 & $\emptyset$ & $\emptyset$ &  R,B \\
    0 & $\emptyset$ & R &  B \\
    0 & $\emptyset$ & B &  R \\
    1 & R & B &  $\emptyset$ \\
    1 & R & $\emptyset$ &  B \\
    1 & B & R &  $\emptyset$ \\
    1 & B & $\emptyset$ &  R \\
    2 & R,B & $\emptyset$ &  $\emptyset$ \\
    \hline
    \end{tabular}
    \end{center}

    \item LHS: A concert has n-distinct people attending. The concert is giving away k-guitars. Of the k-guitars a small subset of size j are signed. Each person can only receive one guitar and each guitar can only be signed once. How many different ways can the unsigned and signed guitars be distributed?
    
    \medskip
    RHS: The same concert has n-distinct people attending. The concert starts by giving away j-guitars that are signed. After the signed guitars are given away, a new subgroup is made of $n-j$ people that didn't receive one. These people are eligible for $k-j$ unsigned guitars. How many ways can the signed and unsigned guitars be distributed.

    \medskip
    For the identity to be valid, $n > k > j$.

    \item The number of size k multisets from an n-set equals the number of size k subsets from an n+k-1 set. The right hand and left hand of the equivalence equation are shown using non-repeating subset notation.
    \begin{equation*}
        \multiset{n}{k} = \binom{n+k-1}{k}, \quad \multiset{k+1}{n-1} = \binom{k+1+n-1-1}{n-1} = \binom{n+k-1}{n-1}
    \end{equation*}

    As such, the multiset equivalence statement can be written as:
    \begin{equation*}
        \binom{n+k-1}{k} = \binom{n+k-1}{n-1}
    \end{equation*}
    
    I need to finish this later and show that with a $[n+k+1]$ set that the complement of the number of $k$ subsets is $n-1$ subsets. So instead, we can count the complement. 
    \end{enumerate}

\begin{problem} % Problem 5
What does $\binom{n-1}{k-1} + \binom{n-2}{k-1} + \binom{n-3}{k-1} + \dots + \binom{k-1}{k-1}$ equal? Make a conjecture and then give a combinatorial proof.
\end{problem}

My conjecture is that:
\begin{equation*}
    \binom{n}{k} = \binom{n-1}{k-1} + \binom{n-2}{k-1} + \binom{n-3}{k-1} + \dots + \binom{k-1}{k-1}
\end{equation*}

\setcounter{problem}{6}
\begin{problem} % Problem 7
Determine the number of solutions to each of the following equations. Assume all $z_i$ are nonnegative integers unless stated otherwise.
\begin{enumerate}[(a)]
    \item $z_1 + z_2 + z_3 + z_4 = 1$
    \item $z_1 + z_2 + 10z_3 = 8$
    \item $z_1 + z_2 + \cdots + z_{20} = 401 \text{ where each } z_i \ge 1$.
\end{enumerate}
\end{problem}

\textbf{Answer:}
\begin{enumerate}[(a)]
    \item $\multiset{4}{1} = 4$
    \item In this problem, $z_3$ can only be 0, if $z_3 > 0$ then $z_1$ or $z_2$ are forced to be negative which violates the constraints. The answer is $\multiset{2}{8} = \binom{9}{8} = 9$
    \item $\multiset{20}{401-20} = \binom{400}{381} =$ big
\end{enumerate}

\section{Chapter 2.3}
\setcounter{problem}{1}
\begin{problem} % Problem 2
For any integer $n \ge 2$, how many onto functions $[n] \rightarrow [n-1]$ are possible? Give a formula that doesn't involve Stirling numbers.
\end{problem}

\textbf{Answer:}
On page 69 of the book, it shows that the special case of Stirling number $S(n, n-1)$ is the following.
\begin{equation*}
    S(n, n-1) = \binom{n}{2}
\end{equation*}
Later, counting onto functions $[k] \rightarrow [n]$ is equal to $S(k,n) \cdot n!$. Putting it all together, the number of onto functions with $[n] \rightarrow [n-1]$ is:

\begin{equation*}
    \binom{n}{2} \cdot (n-1)!
\end{equation*}


\setcounter{problem}{6}
\begin{problem} % Problem 7
Find the number of equivalence relations on an $n$-set.
\end{problem}

\textbf{Answer:}
On page 217 in ``\textit{Book of Proof}'' by Richard Hammack it makes the following statements regarding partitions and equivalence relations. ``\textit{Thus equivalence relations and partitions are really just two different ways of looking at the same thing}''. So really, the question could be rephrased as \textit{Find the number of partitions on an $n$-set}. The number of equivalence relations / partitions is the Bell Number, $B(n)$.
\begin{equation*}
    B(n) = \sum_{k=1}^{n} S(n,k) \quad \text{for all } n \ge 1. 
\end{equation*}

\setcounter{problem}{13}
\begin{problem} % Problem 14
Here is a simple recursive python program for computing $S(n,k)$ based on Theorem 2.3.1. The program assumes $n,k \ge 0$. It works but it is extremely wasteful. Why? Design a more efficient algorithm.
\begin{minted}{python}
def Stirling(n: int, k: int):
    if n == k:
        return 1
    elif n < k:
        return 0
    elif n > 0 and k == 0:
        return 0
    else:
        return Stirling(n-1, k-1) + k*Stirling(n-1, k)
\end{minted}
\end{problem}
This code is poor because a valid value is only returned if $n == k$. Typically this woun't happen,so the recursive function will keep dropping the values for $n$ and $k$ lower until they are equal, $n <k$ or $k==0$. If $n$ is high and $k$ is low, this could take $n-1$ loops to complete.

\medskip
A better solution is to directly calculate the Stirling number with a binomal distribution.
\begin{equation*}
    S(n,k) = \dfrac{1}{k!} \sum_{j=0}^{k} \binom{k}{j} (-1)^j (k-j)^n
\end{equation*}
The pseudo code (no internet to check syntex) with python would be the following.

\begin{minted}{python}
    from itertools import factorial, combination # I think python has built in?
    def Stirling(n: int, k: int):
        bi_sum = 0 # all the stuff
        for j in range(0, k): # check index values
            bi = -1**j*(k-j)**n*combination(k,j)
            bi_sum += bi
        return bi_sum/factorial(k)
\end{minted}

\setcounter{problem}{16}
\begin{problem}
Solve a linear system to find numbers $a, b, c, d, e$ so that the following polynomial equation is true:
\begin{equation*}
    x^4 = a \cdot x_0 + b \cdot x_1 + c \cdot x_2 + d \cdot x_3 + e \cdot x_4
\end{equation*}
In the equation, the following represent $x_1, x_2 \dots x_0$.
\begin{align*}
    x_4 & = x(x-1)(x-2)(x-3) = x^4 - 6x^3 + 11x^2 -6x \\
    x_3 & = x(x-1)(x-2) = x^3 - 3x^2 + 2x\\
    x_2 & = x(x-1) = x^2 - x \\
    x_1 & = x \\
    x_0 & = 1
\end{align*}
Express the solution $a,b,c,d,e$ in terms of numbers studied in this section. Use linear algebra and formulas for expanding the coefficients.
\end{problem}

\textbf{Answer:}

Rewriting the original problem yields.

\begin{equation*}
    x^4 = a \cdot 1 + b \cdot x + c \cdot (x^2 - x) + d \cdot (x^3 - 3x^2 + 2x) + e \cdot (x^4 - 6x^3 + 11x^2 - 6x)
\end{equation*}

Writing the system of equations in matrix form.
\begin{equation*}
    \begin{bmatrix}
        1 & 0 & 0 & 0 & 0 \\
        0 & 1 & -1 & 2 & -6 \\
        0 & 0 & 1 & -3 & 11 \\
        0 & 0 & 0 & 1 & -6 \\
        0 & 0 & 0 & 0 & 1
    \end{bmatrix}
    \begin{bmatrix}
        a \\
        b \\
        c \\
        d \\
        e
    \end{bmatrix}
    =
    \begin{bmatrix}
        0 \\
        0 \\
        0 \\
        0 \\
        1
    \end{bmatrix}
\end{equation*}
Honestly, from here, I would just use back-substitution to solve the linear system. I'm not sure why you need combinatorics or what I am missing. Doing the math in my head...
\begin{align*}
e & = 1 \\
d & = 6 \\
c & = 3*6 - 11 = 7 \\
b & = 7 - 2*6 + 6 = 1 \\
a & = 0
\end{align*}

\end{document}